\subsection{Reproducing the Computation}\label{app:reproduction}

Download the repository archive labeled \emph{Source code (zip)}
on the GitHub release page cited in the paper. The three files named \path{fgst-verification-data-*.zip}
are archived in the
\href{https://github.com/rewere143-cell/fgst-strong-clique/releases/tag/v1.0.0}{version 1.0.0 data release}.
Extract the source archive, open its root
directory, and run \texttt{python restore\_data.py} followed by the
path to the directory containing the downloaded ZIP files. The
programs and their inputs are then in \path{reproducibility/}.
The simple-graph verification programs run offline on a 64-bit
little-endian machine with Python 3.9 or later and a C17 compiler
such as GCC or Clang. From the directory containing
\path{reproducibility/}, run
\begin{verbatim}
cd reproducibility
python run_all.py --output results/fresh-run
\end{verbatim}
with Python assertions enabled. Use an output directory that does
not exist. To select a compiler, pass \verb|--cc| followed by its path.

The runner checks the byte counts and SHA-256 digests against
\path{DATA_SHA256.json}, compiles the C programs and runs the checks
in \textbf{Table~\ref{tab:verification}}. It records the results in
\path{summary.json} in the output directory. A successful run has
\verb|all_passed| set to \verb|true| and \verb|exact_bound| set to
\verb|5/4|. The output directory also contains the individual reports.

\begin{table}[htbp]
\centering
\caption{Verification programs and their checks.}
\label{tab:verification}
\begin{tabular}{@{}lp{0.49\linewidth}@{}}
\toprule
Program & Check\\
\midrule
\path{check_index.c} &
Order-six pattern validity and relabeling orbits\\
\path{check_low_orders_and_bases.py} &
Lower-order indices, degree rows and flag lists\\
\path{check_reverse.c} &
Degree and nonconstant moment coefficients\\
\path{check_certificate.py} &
Exact residual and constant for each rational certificate\\
\path{check_finite_models.py} &
Direct map counts on six finite graphs\\
\bottomrule
\end{tabular}
\end{table}

We verify the index tables before using them to compute
\eqref{eq:integer-residual}. The index programs enumerate all labeled
patterns on three through six positions with pair $01$ marked and
test each pattern for connectivity and admissibility. For a
representative of each stored class, the programs apply all
permutations of the positions, retaining the relabelings with pair
$01$ marked. The programs compare the resulting orbit with the
patterns assigned to the class. The lower-order checker also verifies
that every degree row specifies a connected admissible base pattern
and a position of that pattern. For each listed type and flag order,
the checker verifies that the flag list contains one representative
from every orbit under permutations fixing the type positions.

The program \path{check_reverse.c} starts with the stored
coefficients and subtracts the contributions to them from
\eqref{eq:degree-extension} and \eqref{eq:moment-expansion}. For a
degree row $(F,u)$ with $F$ of order greater than two, it first
adds $1$ to the coefficient of the orbit of $F$, accounting for
$-q_F$ on the left side of the equation. It then enumerates the
connected admissible patterns in the verified tables. A pattern
extending $F$ at $u$ as in \eqref{eq:degree-extension} subtracts
$1$ from the coefficient of its orbit in that row. For a moment
entry, a pattern on three through six positions subtracts $1$
from the coefficient of its orbit if the restrictions to the two
position sets in the proof of \textbf{Lemma~\ref{lem:moment}} are
the prescribed flags. There are $2{,}116{,}783$ such moment
contributions. The program requires every updated key to be in
the table and all remainders to be zero. This checks the
$433{,}899$ stored degree coefficients, the $1{,}895{,}405$ stored
moment coefficients and the zeros at all omitted keys.

The program \path{check_certificate.py} checks the denominators,
the vector lengths and the nonnegativity of the weights. It computes
$B_\ell$ using \eqref{eq:integral-Q}. To reconstruct $o$, the program
counts the labeled patterns of order four with pair $23$ marked and
adds $\frac{1}{2}$ to the corresponding coordinate for each such
pattern, as in \eqref{eq:objective-vector}. It uses integer arithmetic
by scaling these contributions by $\Lambda$, which is even. The
program also accounts for the constant contribution from $E_\star$.
In each block whose type has order two, it includes the first
diagonal entry of $B_\ell$ in \eqref{eq:integer-constant}, since
$(K_\ell)_{11}=1$ and all other entries of $K_\ell$ are zero.
For larger types, $K_\ell=0$.

With the verified coefficients, the program evaluates every entry of
$\rho$ and the constant in \eqref{eq:integer-constant}. Of the
$23{,}789$ entries of $\rho$, $23{,}590$ are positive and $199$ are
zero. The largest is $899{,}382{,}275{,}568$. The constant equals
$3{,}000{,}000{,}000$.

The runner verifies both certificates in \path{data/} separately.
It also checks that the integer verifier rejects two altered versions
of the compact certificate. One has a negative weight numerator.
The other increases the degree multiplier numerator for the
two-position base by one, changing the certified constant.

The checks on finite graphs compare the degree equations, the moment
entries and the objective in \eqref{eq:objective-value} with direct
counts of color-preserving maps. They use $K_2$, $K_3$ and $C_5$ with
all edges marked, and $K_4$ and $K_{3,3}$ with marked matchings.
The final example is an eight-cycle formed by joining corresponding
endpoints of two disjoint four-vertex paths. The two end edges of one
path are marked, and its unmarked middle edge joins their endpoints.

The directories \path{results/reference-run/} and
\path{results/revision-replay/} contain the archived verification reports from two independent runs. The file
\path{research_provenance/README.md} gives commands to regenerate the
coefficient tables and construct a rational certificate from the
included numerical output. That directory also preserves the
original optimization code, its initialization data and the solver
tests. The verification command above uses the fixed tables and
rational certificates in \path{data/} and performs no numerical
optimization.

\subsubsection{Odd-Degree Certificates and Integer Checks}

Extract \path{reproducibility/odd_degree/simple_odd_degree_fixed.zip}
to a fresh directory. From that directory run
\begin{verbatim}
python -B verify_odd_archive.py --output replay
\end{verbatim}
The output directory must not exist. The compiler option \verb|--cc|
is also available for this runner. It checks the hashes of the 69
archived inputs and replays the pattern-index, coefficient and
quadratic-bound checks. It then verifies the five-cycle certificates in
\eqref{eq:odd-moment-bounds} by checking the corresponding inequalities
with rational arithmetic. The program \path{check_c5_dual.py} reads
the verified coefficient tables and the files
\path{c5_abundance_column_candidate_third.json},
\path{c5_abundance_column_candidate_second.json} and
\path{c5_anchor_column_candidate_fifth.json}. The first two bound the
number of marked induced five-cycles from below. The third bounds the
sum of their remaining degrees from above. The file
\path{independent_replay.json} identifies the fixed archive by its
SHA-256 digest.

The runner also executes \path{verify_odd_simple.py}. Its small-degree
calculation calls \path{check_anchored_odd_small.py} with the rational
parameters in \textbf{Appendix~\ref{app:odd-small}}. The four cases
used for \textbf{Lemma~\ref{lem:simple-small-support}} can be replayed
independently.

\noindent\begin{minipage}{\linewidth}
With the fixed archive extracted, execute the following Python statements
from the directory containing the verifiers.
\begin{verbatim}
from fractions import Fraction as F
from check_anchored_odd_small import one_degree
for D in (3, 5, 7, 9):
    failed = one_degree(D, F(99999, 100000), F(56425, 100000),
                        F(2, 10**9), F(1, 10**9),
                        step=1, relax=True)
    assert not failed, (D, failed)
\end{verbatim}
\end{minipage}\par
Here \verb|step=1| checks the upper bound for a counterexample with
$f(D)+1$ edges. The full runner repeats the certificate and finite
checks from the original odd-degree archive. It records
\verb|finite_arithmetic_passed| and \verb|all_passed| on success.

\subsubsection{Multigraph Certificates and Finite Checks}
\label{app:multi-reproduction}

The multigraph verifiers and their inputs are in
\path{reproducibility/multigraph/}.
Its \path{REPLAY.md} indexes the supporting proofs and the fixed
inputs. The source package also includes \path{PROOF_NOTES.md},
which identifies the arguments cited in the appendices. From the
companion directory, with Python assertions enabled, run
\begin{verbatim}
python -X utf8 -B verify.py hashes
python -X utf8 -B verify.py flags
python -X utf8 -B verify.py finite --jobs 2
\end{verbatim}
The first command checks the sizes and SHA-256 digests of the files
listed in the two release manifests. The second rebuilds
the weighted flag calculations for the uniform bound and all eleven
odd-degree intervals in \textbf{Table~\ref{tab:multi-odd}}.
It requires NumPy and SciPy.
It uses the Ore-degree bounds listed in the accompanying proof notes.
Appendix~\ref{app:ore-replay} describes the separate check for the
bipartite bound.
The third command verifies the finite cases used in
\textbf{Appendix~\ref{app:multi-finite}}, including the integer enumeration in
degree five and the positive reverse unit propagation proofs for the
Boolean formulas. The proof checker verifies every added clause by
unit propagation and checks that the empty clause is derived. The
source package gives the archive digest and the independently replayed
checks in \path{VERIFICATION.md}.

In the same directory, the following command verifies
\eqref{eq:absence-red-cycle}. It reconstructs the constraints for
$q_C=0$, computes the rational upper bound on $M$ from the stored
coefficients and checks that it is strictly less than $\frac{98}{81}$.
\begin{verbatim}
python -X utf8 -B check_induced_five_cycle.py
\end{verbatim}

\subsubsection{Rational Interval Check for the Bipartite Ore Bound}
\label{app:ore-replay}

The bound with coefficient $\frac{2973}{10000}$ uses induction on the
number of marked edges in a loopless multigraph whose marked graph is
bipartite. The companion
\href{https://github.com/rewere143-cell/fgst-strong-clique/blob/v1.0.0/reproducibility/multigraph/\%E4\%B8\%80\%E8\%88\%AC\%E6\%97\%A0\%E7\%8E\%AF\%E5\%A4\%9A\%E9\%87\%8D\%E5\%9B\%BE\%E7\%9A\%84\%E6\%8E\%A8\%E5\%B9\%BF/\%E4\%BA\%8C\%E9\%83\%A8\%E6\%A0\%87\%E8\%AE\%B0\%E8\%A7\%81\%E8\%AF\%81\%E4\%BA\%8F\%E6\%8D\%9F\%E4\%B8\%8E\%E9\%82\%BB\%E5\%9F\%9F\%E5\%B7\%AE.md}{proof on neighborhood differences}
derives the inequalities for the interval check. Equations~(1)--(6)
of that proof give the count of edges joining endpoints of disjoint
marked edges and the two estimates for the number of adjacent pairs
in the root's part. The penultimate section derives the one-sided
weighted inequalities used by the option \verb|--tight|. The last
section gives the commands for the coefficient $\frac{2973}{10000}$.
In the
multigraph proof directory, \path{check_multigraph_bipartite_coupled.py}
checks the initial coefficient $\frac{2989}{10000}$.
\path{check_multigraph_recursive_ore.py} then verifies the coefficient
$\frac{1489}{5000}$ and writes a table of conditional bounds.
\path{refine_recursive_ore_child_ratio.py} refines this table in four
rounds, with each round using the bounds from the preceding round.
For the final run, \path{check_bipartite_neighborhood_slack.py} takes
the refined table through \verb|--conditional-table| and uses
\verb|--prior-beta 1489/5000|, \verb|--beta 2973/10000|,
\verb|--parts 4| and \verb|--tight|. All interval comparisons use
rational endpoints. The final check completes after $73{,}919$
subdivisions with no unresolved boxes. The output is in
\path{verification/ore_final.log}.

The coefficient $\frac{273323}{10^6}$ for the base copies comes
from a separate certificate for a simple host with bipartite marked
edge set. In the simple-graph snapshot indexed by
\path{reproducibility/multigraph/REPLAY.md}, open
\path{source_snapshot/ore_simple/}. The file
\path{ore_certificate_proof.md} gives the normalization, the
marked-edge degree inequalities and the exact dual identity.
From that directory run \texttt{python -B verify\_ore\_certificate.py}
with two arguments, the input \path{ore_plain_repaired_certificate.json}
and a new output path, in that order.
The certified system has $12{,}772$ coordinates, $3{,}120$ inequalities
and $47$ moment matrices. The reference report
\path{verification/simple_ore.json} records the checks on the
nonnegative multipliers and positive semidefinite factors, and gives
the corrected rational bound.
